% ==========================================================
\section{问题 3：初糖浓度的影响规律与工艺优化}

\subsection{终点产量随初糖的变化}

三个候选模型（表 \ref{tab:q3end}）：线性模型假设“糖多酸多、永不饱和”；Michaelis--Menten 型假设“逐渐饱和”；“产率—上限”模型则假设“在效价上限以下按固定得率转化，超过上限后被抑制卡住”。

\begin{table}[H]
\centering
\caption{终点产量模型的比较（$n=6$）}
\label{tab:q3end}
\small
\begin{tabular}{llcccc}
\toprule
编号 & 模型 & 参数估计 & $R^2$ & RMSE & AIC \\
\midrule
C1 & $P=aS_0$ & $a=0.846$ & 0.9027 & 9.86 & 29.46 \\
C2 & $P=\dfrac{aS_0}{b+S_0}$ & $a=299.4,\ b=241.5$ & 0.9763 & 4.87 & 22.99 \\
C3 & $P=\min\{YS_0,\ P_{cap}\}$ & $Y=0.9614,\ P_{cap}=109.9$ & \textbf{0.99970} & \textbf{0.54} & $\mathbf{-3.29}$ \\
\bottomrule
\end{tabular}
\end{table}

C1 在 150 g/L 处预测 $126.9$ g/L，比实测高 $17$ g/L，说明“线性无上限”不合格；C2 形式光滑但系统性高估低糖端、低估高糖端；C3 只用两个参数就同时刻画了两段行为，RMSE 仅 $0.54$ g/L。

\begin{equation}
P_\infty(S_0)=\min\{0.9614\,S_0,\ 109.9\}.
\label{eq:pend}
\end{equation}

由此得到一个重要的“平衡点”：
\begin{equation}
S_0^*=\frac{P_{cap}}{Y}=\frac{109.9}{0.9614}=114.3\ \mathrm{g/L}.
\label{eq:s0star}
\end{equation}
当 $S_0<S_0^*$ 时糖被用完（$S_r\approx0$），产量与初糖成正比；当 $S_0>S_0^*$ 时乳酸先撞上 $109.9$ g/L 的天花板，剩余糖无法转化。式 \eqref{eq:pend} 预测 150 g/L 时残糖
$$S_r=S_0-\frac{P_{cap}}{Y}=150-114.3=35.7\ \mathrm{g/L},$$
与实测 $35.7$ g/L \wj{完全一致}（图 \ref{fig:q3}(a)(d)）。

\subsection{速率随初糖的变化}

\paragraph{(1) 最大产酸速率。}最大速率对初糖不敏感但高糖端下降（$3.2\sim3.4\to2.6$），这是典型的\wj{底物抑制}（高渗透压、高糖对转运蛋白的抑制）。用 Haldane 方程
\begin{equation}
r_{max}(S_0)=\frac{V_{max}S_0}{K_s+S_0+S_0^2/K_i}
\label{eq:haldane}
\end{equation}
拟合得 $V_{max}=4.66$、$K_s=7.68$、$K_i=224.9\ \mathrm{g/L}$，$R^2=0.769$。$K_i$ 的存在表明超过约 $150$ g/L 后速率将急剧下降，这也解释了为什么 150 g/L 那一批的 $r_{max}$ 明显偏低。

\paragraph{(2) 平均产酸速率为何下降。}平均速率 $\bar r=P_\infty/t_e$ 单调下降（$2.18\to1.02$），其机理并非“菌越来越弱”，而是\wj{时间被“尾巴”拉长}：产酸速率随 $P$ 接近 $P_{cap}$ 而按 $(1-P/P_{cap})$ 衰减，越接近上限越慢，平均下来必然下降。下面用统一的动力学模型定量说明。

\subsection{时长与速率的统一动力学模型}
把“底物限制 + 底物抑制 + 产物抑制”写进一个速率表达式：
\begin{equation}
\frac{\dd P}{\dd t}=\underbrace{\frac{V_{max}S}{K_s+S+S^2/K_i}}_{\text{底物项}}\cdot\underbrace{\Big(1-\frac{P}{P_{cap}}\Big)}_{\text{产物抑制}},\qquad S=S_0-\frac{P}{Y}.
\label{eq:q3ode}
\end{equation}
“发酵终点”按试验二的判定标准（残糖基本不再下降）定义为速率降到阈值 $v_{stop}$ 的时刻：
\begin{equation}
t_e=\inf\Big\{t:\ \frac{\dd P}{\dd t}<v_{stop}\Big\},\qquad P_\infty=P(t_e),\qquad \bar r=\frac{P_\infty}{t_e}.
\label{eq:te}
\end{equation}
给定 $\boldsymbol\theta=(V_{max},K_s,K_i,P_{cap},Y,v_{stop})$，先由 $\dd P/\dd t=v_{stop}$ 解出终点浓度 $P_\infty$（二分法），再由 $t_e=\int_0^{P_\infty}\big(\dd P/\dd t\big)^{-1}\dd P$ 积分得到时长。这样定义使 $t_e$ 对参数连续可微，便于最小二乘收敛。

联合拟合 6 个初糖水平下的终点时间、终点产量与最大速率（相对误差加权），得表 \ref{tab:q3kin}。

\begin{table}[H]
\centering
\caption{问题 3 动力学模型参数}
\label{tab:q3kin}
\small
\begin{tabular}{lccc}
\toprule
参数 & 估计值 & 标准误 & 含义 \\
\midrule
$V_{max}$ & 3.657 g/(L$\cdot$h) & 0.441 & 最大比产酸速率 \\
$K_s$     & 2.283 g/L & 1.116 & 半饱和常数 \\
$K_i$     & 513.6 g/L & 397.9 & 底物抑制常数 \\
$P_{cap}$ & 111.55 g/L & 3.04 & 产物抑制上限 \\
$Y$       & 0.983 g/g & 0.017 & 乳酸得率 \\
$v_{stop}$& 0.1239 g/(L$\cdot$h) & 0.044 & 终点速率阈值 \\
\bottomrule
\end{tabular}
\end{table}

\begin{table}[H]
\centering
\caption{动力学模型对试验二的拟合（括号内为实测值）}
\label{tab:q3pred}
\small
\begin{tabular}{ccccccc}
\toprule
$S_0$ & $t_e$/h & $P_\infty$/(g/L) & $r_{max}$ & $\bar r$ & $S_r$/(g/L) \\
\midrule
20  & 9.8 (9)    & 19.56 (19.6) & 3.17 (3.20) & 1.993 (2.180) & 0.10 \\
40  & 18.7 (18)  & 39.20 (37.3) & 3.22 (3.30) & 2.091 (2.070) & 0.13 \\
60  & 30.7 (33)  & 58.81 (57.9) & 3.17 (3.40) & 1.918 (1.750) & 0.18 \\
80  & 48.5 (63)  & 78.36 (77.4) & 3.09 (3.10) & 1.617 (1.230) & 0.29 \\
100 & 81.4 (90)  & 97.48 (96.0) & 3.00 (3.30) & 1.197 (1.070) & 0.84 \\
150 & 116.7 (108)& 107.26 (109.9)& 2.80 (2.60) & 0.919 (1.020) & 40.90 \\
\midrule
$R^2$ & 0.954 & 0.998 & 0.546 & 0.819 & -- \\
RMSE  & 7.82 h & 1.56 g/L & 0.177 & 0.200 & -- \\
\bottomrule
\end{tabular}
\end{table}

\paragraph{评价。}模型对终点产量与平均速率的刻画很好（$R^2=0.998$ 与 $0.819$），对终点时间的 $R^2=0.954$。最大的偏差出现在 $S_0=80$ g/L（预测 $48.5$ h，实测 $63$ h，低 $23\%$）与 $S_0=20$ g/L（预测 $9.8$ h，实测 $9$ h）。对 80 g/L 这一点，相邻两档的时间分别为 33 h 与 90 h，而 80 g/L 的 63 h 明显“跳档”，很可能是该批试验的操作差异（如接种量、溶氧或补碱记录）所致；由于只有一个重复，本文不为其增加额外的经验修正项，而是在优化中做稳健性检验（见 \ref{sec:robust}）。\wj{这一点也说明：6 个数据点拟合 6 个参数，模型对单点异常是敏感的，报告残差比报告 $R^2$ 更重要。}

\subsection{高糖残糖明显剩余的原因}
综合以上三方面，150 g/L 时残糖 $35.7$ g/L 是三重因素叠加的结果，但主次分明：

\begin{enumerate}[leftmargin=2.2em,itemsep=0.12em]
\item \wj{主因——产物抑制形成的效价天花板。}由式 \eqref{eq:pend}，乳酸一旦达到 $109.9$ g/L，$(1-P/P_{cap})\to0$，无论还有多少糖都无法继续转化。$150-114.3=35.7$ g/L 的糖被“剩”下，数量与被抑制“卡住”的部分严格对应。这是数量上的确证，而不是定性猜测。
\item \wj{次因——底物（渗透）抑制拖慢进程。}$K_i=514$ g/L 意味着一开始速率就比低糖时低约 $12\%$，使到达天花板需要更长时间（$116.7$ h vs 90 h），给“酸累积到上限”创造了条件。
\item \wj{诱因——高糖对菌体得率的损耗。}OD$/S_c$ 由 $0.10$ 降到 $0.085$，说明一部分糖与能量被用于对抗渗透胁迫（相容性溶质合成、质子泵耗能），而非用于产酸。
\end{enumerate}

需要强调：\wj{残糖剩余不是因为“菌吃不动糖”，而是因为“酸到了上限、菌被自己的产物抑制”}。这一点可由“150 g/L 的最大产酸速率仍有 2.6，只是持续时间被拉长”得到侧面印证——菌并非失去活力，而是产出被抑制项“掐断”。

\subsection{工艺优化}

\paragraph{(1) 目标函数。}企业关心的是“单位时间赚多少”，而不是单纯的产量或速率。设乳酸售价 $p_P=0.010$ 元/g（10 元/kg）、葡萄糖 $p_S=0.0032$ 元/g（3.2 元/kg）、下游按产品计的变动成本 $c_{var}=0.0015$ 元/g、按发酵液体积计的处理成本 $c_{vol}=0.12$ 元/L（含蒸发、离交、结晶的固定开销，体现“效价越低、每千克产品要处理的水越多”）、发酵运行成本 $c_{op}=0.0015$ 元/(L$\cdot$h)、批次辅助时间 $t_0=6$ h（清洗、灭菌、接种、放罐）。则以“每升发酵液每小时的净收益”为目标：
\begin{equation}
\max_{S_0,\,t}\ J(S_0,t)=\frac{p_P P(t)-p_S S_0-c_{var}P(t)-c_{vol}-c_{op}\,t}{t+t_0},
\label{eq:J}
\end{equation}
其中 $P(t)$ 由式 \eqref{eq:q3ode} 给出。同时考察“单位产品成本”
\begin{equation}
C(S_0)=\frac{p_S S_0+c_{op}t_e+c_{var}P_\infty+c_{vol}}{P_\infty}\times1000\quad[\text{元/kg}].
\label{eq:C}
\end{equation}

\paragraph{(2) 结果。}在 $S_0\in[20,150]$、$t\in[2,t_e]$ 的网格上寻优，得到（图 \ref{fig:q3opt}）：

\begin{table}[H]
\centering
\caption{不同目标下的最优工艺参数}
\label{tab:opt}
\small
\setlength{\tabcolsep}{3pt}
\begin{tabular}{llccc}
\toprule
优化目标 & 含义 & 最优 $S_0$/(g/L) & 最优 $t$/h & 指标值 \\
\midrule
$\max J$ & 单位时间净收益 & 74 & 36.5 & 0.00440 元/(L$\cdot$h) \\
$\min C$ & 单位产品成本 & 89 & 60.3 & 7.19 元/kg \\
$\max \bar r$ & 生产强度 & 66 & 18 & 1.88 g/(L$\cdot$h) \\
$\max P_\infty$ & 终点产量 & 143 & 115.9 & 107.3 g/L \\
\bottomrule
\end{tabular}
\end{table}

\textbf{推荐方案}：\wj{初糖 $S_0\approx 75\sim90$ g/L，放罐时间 $40\sim60$ h}（以残糖降到 1 g/L 以下或产酸速率降到 $0.2\ \mathrm{g/(L\cdot h)}$ 为放罐判据）。这一区间同时覆盖了“单位时间收益最大”（$74$ g/L、$36.5$ h）与“单位产品成本最低”（$89$ g/L、$60.3$ h）两个最优点。理由如下：

\begin{enumerate}[leftmargin=2.2em,itemsep=0.10em]
\item \wj{不浪费糖。}$75\sim90$ g/L 低于平衡点 $114.3$ g/L，糖可全部转化，得率 $0.96$ 以上；一旦超过 $114.3$，每多加 1 g 糖只是把 1 g 糖原样留在罐里，还要多付糖钱与后续分离成本。
\item \wj{效价足够高。}$75\sim90$ g/L 初糖对应终点效价 $73\sim87$ g/L，已达到下游经济蒸发/结晶通常要求的 $70\sim80$ g/L 门槛，$c_{vol}$ 的稀释惩罚退居次要。
\item \wj{时间成本合适。}继续提高初糖到 $110$ g/L，效价只增加 $20\%$，时间却增加 $70\%$（$60\to100$ h），设备周转显著恶化；而降低到 $60$ g/L，时间虽短但单位产品成本反而上升（$c_{vol}$ 摊薄不足）。
\item \wj{抗风险。}该区间处于 $J$ 平台区（$70\sim95$ g/L 内 $J$ 只变化 $11\%$，远小于两端的变化幅度），即使价格、成本估计有偏差，结论也不会翻转。
\end{enumerate}

\begin{figure}[H]
\centering
\includegraphics[width=0.99\textwidth]{fig5_q3_opt.png}
\caption{问题 3 的优化结果。(a) 产量—强度权衡（Pareto 曲线）；(b) 单位产品成本与单位时间收益；(c) 最优初糖对成本参数的敏感性；(d) 初糖—放罐时间的二维寻优}
\label{fig:q3opt}
\end{figure}

\paragraph{(3) 敏感性。}改变产品价格 $p_P$ 与下游体积成本 $c_{vol}$（其余参数不变），最优初糖在 $58\sim109$ g/L 之间移动（图 \ref{fig:q3opt}(c)）：下游成本越高（$c_{vol}$ 大），越应提高效价，最优初糖右移；产品价格越高，越值得“多花时间多出活”，最优初糖也略右移。但\wj{“不超过 114 g/L”这一硬边界始终成立}——这是模型给出的、不依赖任何价格假设的结构性结论。

\subsection{推荐方案的稳健性}
\label{sec:robust}
为检验结论是否被 80 g/L 那个异常点带偏，做留一重拟合：依次去掉一个初糖水平后重新估计全部 6 个参数，再重新寻优，结果见表 \ref{tab:loo}。

\begin{table}[H]
\centering
\caption{留一重拟合的稳健性（每列去掉一个初糖水平后重新建模）}
\label{tab:loo}
\small
\begin{tabular}{lccccccc}
\toprule
去掉的 $S_0$ & 20 & 40 & 60 & 80 & 100 & 150 & 全样本 \\
\midrule
最优 $S_0$（$\max J$） & 76 & 74 & 73 & 73 & 75 & 69 & 73 \\
最优 $S_0$（$\min C$） & 90 & 90 & 91 & 89 & 89 & 81 & 90 \\
\bottomrule
\end{tabular}
\end{table}

可见：去掉任何一个点，收益最优初糖都落在 $69\sim76$ g/L，成本最优初糖都落在 $81\sim91$ g/L，推荐区间 $75\sim90$ g/L 始终被覆盖。\wj{推荐结论对单个数据点不敏感}，即使 80 g/L 那一批确为异常，也不改变建议。

\begin{figure}[H]
\centering
\includegraphics[width=0.98\textwidth]{fig4_q3_endpoints.png}
\caption{问题 3 的终点规律。(a) 终点产量；(b) 终点时间；(c) 平均产酸速率；(d) 终点残糖}
\label{fig:q3}
\end{figure}

% ==========================================================
\section{问题 4：过程强化策略——及时移出乳酸}

\subsection{为什么要选“移出乳酸”}
问题 2 与问题 3 指向同一个结论：\wj{两类菌株的瓶颈都是产物抑制}。试验一中乳酸达到 $9.9$ g/L 就基本停产（残糖 $25$ g/L）；试验二中乳酸达到 $110$ g/L 就封顶（高糖时残糖 $35.7$ g/L）。“补加糖料”针对的是“糖不够”，而本问题恰恰\wj{不缺糖、缺的是“容纳酸的能力”}，因此补料只能避免浪费、不能突破天花板；只有把酸从罐内移走，才能直接解除限制。

本文选择\wj{原位产物分离（ISPR，及时移出部分乳酸）}作为强化策略：在发酵过程中用树脂吸附、电渗析或溶剂萃取把乳酸连续移出（或部分移出后中和），使罐内乳酸浓度维持在一个较低水平。

\subsection{模型}
在问题 2 的模型上增加一个“移出项”。设罐内乳酸被控制在 $P_{sp}$：
\begin{equation}
\frac{\dd P}{\dd t}=\Big(\alpha\frac{\dd L}{\dd t}+\beta X\Big)\Big(1-\frac{P}{P_m}\Big)-u(t),\qquad
u(t)=\max\Big\{0,\ \frac{\dd P}{\dd t}\Big|_{u=0}\Big\}
\ \text{当}\ P\ge P_{sp}.
\label{eq:ispr}
\end{equation}
实践中即“产出多少移走多少，维持 $P\le P_{sp}$”。累计产量为 $\Pi(t)=P_{sp}+\int_0^t u\,\dd s$，耗糖仍满足 $\dd S/\dd t=-\dfrac{1}{Y_{P/S}}\dfrac{\dd P}{\dd t}-m_s$，当 $S$ 降到 $0.5$ g/L 时停止。用四阶定步长积分（$\Delta t=0.005$ h）求解。

\subsection{结果}
\begin{table}[H]
\centering
\caption{不同乳酸控制水平下的单批效果（试验一条件）}
\label{tab:ispr}
\small
\setlength{\tabcolsep}{5pt}
\begin{tabular}{ccccccc}
\toprule
$P_{sp}$/(g/L) & 累计产酸/(g/L) & 移出量/(g/L) & $t_e$/h & $S_r$/(g/L) & $\bar r$/(g$\cdot$L$^{-1}\cdot$h$^{-1}$) & 倍数 \\
\midrule
不控（基线） & 9.72 & 0 & 14.0 & 25.1 & 0.694 & 1.00 \\
2 & 35.70 & 33.70 & 14.9 & 0.50 & 2.393 & 3.67 \\
4 & 35.59 & 31.59 & 17.2 & 0.49 & 2.067 & 3.66 \\
6 & 35.38 & 29.38 & 21.5 & 0.50 & 1.644 & 3.64 \\
8 & 34.80 & 26.80 & 33.5 & 0.50 & 1.037 & 3.58 \\
9 & 33.68 & 24.68 & 56.5 & 0.50 & 0.596 & 3.47 \\
\bottomrule
\end{tabular}
\end{table}

结论非常清晰（图 \ref{fig:q4}(a)(b)）：

\begin{enumerate}[leftmargin=2.2em,itemsep=0.10em]
\item \wj{单批产量提高到 3.6 倍。}从 $9.72$ 提高到 $35.6\sim35.7$ g/L，接近“糖全部转化”的理论上限 $Y_{P/S}S_0\approx36.9$ g/L；残余糖从 $25.1$ 降到 $0.5$ g/L，糖的利用率从 $28\%$ 提到 $98\%$。
\item \wj{产出强度提高到 3 倍。}控制在 $4$ g/L 时，平均产出强度从 $0.694$ 提高到 $2.07\ \mathrm{g/(L\cdot h)}$，即同样的罐、同样的时间，产量翻三倍。
\item \wj{存在经济最优的控制水平。}控制得越严（$P_{sp}$ 越小），抑制解除越彻底、产量略高，但需要处理的酸量也越大、时间并不一定最短；$P_{sp}=2\sim6$ g/L 是“高产出—可承受负荷”的折中区，$P_{sp}=8\sim9$ g/L 则收益明显衰减（$1.04$ 与 $0.60$）。
\end{enumerate}

\begin{figure}[H]
\centering
\includegraphics[width=0.99\textwidth]{fig6_q4.png}
\caption{问题 4 的结果。(a) 分批与移酸的累计产酸曲线；(b) 控制水平与产出强度的关系；(c) 补料分批的轨迹}
\label{fig:q4}
\end{figure}

\subsection{对照：补料分批}
作为另一种策略，补料分批（初始 $40$ g/L，按“维持残糖 $30$ g/L”反馈补入浓糖液，总糖 $115$ g）在问题 3 的模型下模拟：终点体积 $1.150$ L、乳酸 $97.5$ g/L、残糖 $0.84$ g/L、历时 $83.9$ h，总产物 $112.1$ g（图 \ref{fig:q4}(c)）。

与“一次投 150 g/L”相比：总产物 $112.1$ g vs $107.3$ g，时间 $83.9$ h vs $115.9$ h，容积产率 $1.34$ vs $0.93\ \mathrm{g/(L\cdot h)}$（提高 $44\%$），且完全避免了 $35.7$ g/L 的糖浪费。但与移酸相比，补料分批\wj{无法突破 $P_{cap}=111.5$ g/L 的天花板}，产物总量仍被“每升最多装 110 g 酸”限制，其增益（$+44\%$）明显小于移酸（$+200\%$ 以上）。

\subsection{工程可行性与成本粗算}
以电渗析为例，从发酵液中分离乳酸的单位能耗约 $0.5\sim1.0$ kWh/kg 酸。按 $0.6$ 元/kWh 计，每千克乳酸增加 $0.3\sim0.6$ 元成本，相对 $10$ 元/kg 的售价仅占 $3\%\sim6\%$；而产量提高到 $3.6$ 倍所摊薄的设备折旧与人工远大于此。即便考虑树脂再生、膜更换与酸损失（按 $5\%$ 计），该策略在经济上仍然成立。\wj{这也说明：当瓶颈是产物抑制时，“把产物拿走”通常比“再优化投料”更有效。}

\subsection{模型的局限与适用范围}

\begin{enumerate}[leftmargin=2.2em,itemsep=0.12em]
\item \wj{菌体子模型是“外生”的。}问题 2 中 $L(t)$ 由试验一直接给出，没有考虑“移除乳酸后菌体会继续生长”。真实的 ISPR 过程中，抑制解除后菌体还会二次生长，产酸能力会进一步提高。因此本文给出的 $35.7$ g/L 是\wj{偏保守的下界}（按“菌体量不增”计算）。
\item \wj{抑制形式是经验假设。}$(1-P/P_m)$ 只是众多可能形式之一（也可用 $K_i/(K_i+P)$ 或指数衰减）。不同形式对“接近上限”的行为影响较大，本文用残差与自助法界定了其不确定性（$P_m\in[9.42,10.48]$）。
\item \wj{两套数据不可混用。}试验一（L-1、不控 pH）与试验二（L-2、pH 9.0）的菌株与控制方式都不同，本文始终分开建模，只在“机理结构”层面互相印证，不交换参数。
\item \wj{未考虑流变与传质。}高浓度发酵液黏度上升会降低传质与搅拌效率，模型中的 $V_{max}$ 实际上是“表观”值，会随规模放大而下降。
\item \wj{数据量小。}试验一只有 8 个点、试验二只有 6 个水平且无重复，参数置信区间普遍较宽（如 $K_i$ 的相对误差达 $77\%$），任何超出试验范围的外推都需谨慎。
\end{enumerate}

\noindent\textbf{适用范围}：本文模型适用于“单一碳源、产物抑制主导、恒温、pH 受控或弱控”的分批/补料分批乳酸发酵，可用于工艺参数的初步筛选与趋势判断；不适用于连续发酵、多底物共代谢或产物毒性极强的极端工况。
